\section{Heisenberg generators}\label{section4}

\subsection[Affine Lie algebra $\hat{\gl}_N$]{Affine Lie algebra $\boldsymbol{\hat{\gl}_N}$}
Let $\gl_N$ be the complex general linear Lie algebra consisting of $N \times N$ matrices. We denote by $E_{i,j}$ the matrix unit with $(i,j)$-th entry~$1$, and set $\bfid = \sum\limits_{i=1}^N E_{i,i}$. The indices~$i$,~$j$ of $E_{i,j}$ are regarded as elements of $\bbZ/N\bbZ$. The transpose of an element $X$ of $\gl_N$ is denoted by ${}^t X$.

Let $\hat{\gl}_N = \gl_N \otimes \bbC[t,t^{-1}] \oplus \bbC c$ be the affine Lie algebra whose Lie bracket is given by
\begin{gather*}
\big[X \otimes t^{r}, Y \otimes t^{s}\big] = [X,Y] \otimes t^{r+s} + r \delta_{r+s,0} \tr(XY) c, \qquad \text{$c$ is central}.
\end{gather*}
We denote the element $X \otimes t^s$ by $X(s)$. We identify the generators as usual:
\begin{alignat*}{4}
& x_0^+= E_{N,1}(1), \qquad && x_0^{-} = E_{1,N}(-1), \qquad && h_0= E_{N,N} - E_{1,1} + c,&\\
& x_{i}^{+}= E_{i,i+1}, \qquad && x_i^{-}= E_{i+1,i}, \qquad &&h_i= E_{i,i}-E_{i+1,i+1}, \qquad i \neq 0.&
\end{alignat*}

We define automorphisms analogous to $\omega$ and $\rho$ for $\hat{\gl}_N$. Let $\omega$ be the anti-automorphism of~${U}\big(\hat{\gl}_N\big)$ defined by $\omega(X(s))= {}^{t}X(-s)$ and $\omega(c)=c$.
The assignment
\begin{gather*}
x_i^{\pm} \mapsto x_{i-1}^{\pm}, \qquad h_i \mapsto h_{i-1}, \qquad c \mapsto c, \qquad \bfid(s) \mapsto \bfid(s) + \delta_{s,0}c
\end{gather*}
gives an algebra automorphism $\rho$ of $U\big(\hat{\gl}_N\big)$.

\begin{lem}\label{lem:rho}We have $\rho(E_{ij}(s)) = E_{i-1,j-1}(s+\delta_{i,1}-\delta_{j,1}) + \delta_{s,0}\delta_{i,1}\delta_{j,1}c$.
\end{lem}
\begin{proof}We can show the identities for $i \neq j$ inductively from those for the Chevalley generators. Then we can show
\begin{gather}
\rho(h_i(s)) = \begin{cases}
h_{-\theta}(s) & \text{if $i=1$},\\
h_{i-1}(s) & \text{otherwise}
\end{cases}\label{eq:rho_h}
\end{gather}
for $s \neq 0$. Indeed we have $\rho(h_1(s)) = \rho([E_{1,2}(s), E_{2,1}]) = [E_{N,1}(s+1), E_{1,N}(-1)] = h_{-\theta}(s)$ for $i=1$. The other cases are similarly proved. The identity (\ref{eq:rho_h}) will be used later.

Let us consider the case $i = j$. The case $s=0$ is proved as follows. Note that the identity $\bfid = \sum\limits_{i=1}^{N-1} i h_i + N E_{N,N}$ holds. Applying $\rho$ to the both sides, we obtain
\begin{gather*}
\bfid + c = \sum_{i=1}^{N-1} i h_{i-1} + N \rho(E_{N,N}).
\end{gather*}
The right-hand side is equal to
\begin{gather*}
\sum_{i=1}^{N-1} i h_i - N h_{N-1} + c + N \rho(E_{N,N}),
\end{gather*}
hence we have $\rho(E_{N,N}) = E_{N,N} + h_{N-1} = E_{N-1,N-1}$. Then we can inductively show
\begin{gather*}
\rho(E_{i,i}) = \rho(h_i + E_{i+1,i+1}) = h_{i-1} + E_{i,i} = E_{i-1,i-1}
\end{gather*}
for $i=N-1,N-2,\ldots,2$. For $i=1$, we have $\rho(E_{1,1}) = \rho(h_1 + E_{2,2}) = h_{0} + E_{1,1} = E_{N,N} + c$. The case $s \neq 0$ is similarly proved by considering
\begin{gather*}
\bfid(s) = \sum_{i=1}^{N-1} i h_i(s) + N E_{N,N}(s)
\end{gather*}
and (\ref{eq:rho_h}).
\end{proof}

We similarly define an algebra automorphism
\begin{gather*}
T_i = \exp\ad x_i^+ \exp\ad\big({-}x_i^-\big) \exp\ad x_i^+
\end{gather*}
of $U\big(\hat{\gl}_N\big)$.

\begin{lem}\label{lem:higher}We have $T_i(X(s))=T_i(X)(s)$ for $X \in \gl_N$ if $i \neq 0$.
\end{lem}
\begin{proof}Obvious from the definition of $T_i$ and the Lie bracket of $\hat{\gl}_N$.
\end{proof}

\begin{lem}\label{lem:diagonal}We have
\begin{gather*}
T_{0}(E_{j,j}(s)) = \begin{cases}
E_{1,1}(s)-\delta_{s,0}c &\text{if } j=N, \\
E_{N,N}(s)+\delta_{s,0}c &\text{if } j=1, \\
E_{j,j}(s) &\text{otherwise}
\end{cases}
\end{gather*}
and
\begin{gather*}
T_{i}(E_{j,j}(s)) = \begin{cases}
E_{i+1,i+1}(s) &\text{if } j=i, \\
E_{i,i}(s) &\text{if } j=i+1,\\
E_{j,j}(s) \ &\text{otherwise},
\end{cases}\qquad i \neq 0.
\end{gather*}
\end{lem}
\begin{proof}We show the case $i \neq 0$. By Lemma~\ref{lem:higher}, it is enough to consider the case $s=0$. Apply~$T_i$ to the identity $\bfid = \sum\limits_{j=1}^{N-1} j h_j + N E_{N,N}$. Then the left-hand side is $T_i(\bfid) = \bfid$ and the right-hand side is
\begin{gather*}
T_i \left(\sum_{j=1}^{N-1} j h_j + N E_{N,N}\right) = \begin{cases}
\displaystyle\sum_{j=1}^{N-1} j h_j - N h_{N-1} + N T_{N-1}(E_{N,N}) & \text{if $i = N-1$}, \vspace{1mm}\\
\displaystyle\sum_{j=1}^{N-1} j h_j + N T_i(E_{N,N}) & \text{otherwise}.
\end{cases}
\end{gather*}
This shows
\begin{gather*}
T_i (E_{N,N}) = \begin{cases}
E_{N-1,N-1} & \text{if $i = N-1$}, \\
E_{N,N} & \text{otherwise}.
\end{cases}
\end{gather*}
Now let $i=1$. We can inductively show that $E_{N-1,N-1}, E_{N-2,N-2}, \ldots, E_{3,3}$ are invariant under~$T_1$ and
\begin{gather*}
T_1(E_{2,2}) = T_1(h_2 + E_{3,3}) = h_2 + h_1 + E_{3,3} = E_{1,1},\\
T_1(E_{1,1}) = T_1(h_1 + E_{2,2}) = -h_1 + E_{1,1} = E_{2,2}.
\end{gather*}
Thus we have shown the assertion for $i=1$. Similarly we can show the other cases.

The case $i=0$ is obtained from the case $i=1$ by applying $\rho$ and using Lemma~\ref{lem:rho}.
\end{proof}

\begin{lem}\label{lem:0_to_k}We have $T_0 T_1 \cdots T_{k-1} (x_k^+) = \begin{cases}
E_{N,k+1}(1) & \text{if $0 \leq k \leq N-2$},\\
-E_{N,1}(2) & \text{if $k=N-1$}.
\end{cases}$
\end{lem}
\begin{proof} First we prove the assertion for $0 \leq k \leq N-2$ by induction on $k$. The case $k=0$ is trivial. Assume that it holds for~$k$, then
\begin{gather*}
T_0 T_1 \cdots T_{k-1} T_{k} \big(x_{k+1}^+\big) = T_0 T_1 \cdots T_{k-1} \big(\big[x_{k}^+, x_{k+1}^+\big]\big) = \big[T_0 T_1 \cdots T_{k-1}\big(x_{k}^+\big), x_{k+1}^+\big] \\
\hphantom{T_0 T_1 \cdots T_{k-1} T_{k} \big(x_{k+1}^+\big)}{} = [E_{N,k+1}(1), E_{k+1,k+2}] = E_{N,k+2}(1).
\end{gather*}
Next we prove the assertion for $k=N-1$. We have
\begin{gather*}
T_0 T_1 \cdots T_{N-2} \big(x_{N-1}^+\big) = T_0 T_1 \cdots T_{N-3} \big(\big[x_{N-2}^+, x_{N-1}^+\big]\big) = \big[T_0 T_1 \cdots T_{N-3}\big(x_{N-2}^+\big), T_0\big(x_{N-1}^+\big)\big] \\
\hphantom{T_0 T_1 \cdots T_{N-2} \big(x_{N-1}^+\big)}{} = [E_{N,N-1}(1), -E_{N-1,1}(1)] = -E_{N,1}(2).
\end{gather*}
The proof is complete.
\end{proof}

\begin{lem}Let $i \leq j$. We have
\begin{gather}
T_{i} T_{i+1} \cdots T_{j-1} \big(x_{j}^+\big) = E_{i,j+1},\qquad T_{i} T_{i+1} \cdots T_{j-1} \big(x_{j}^-\big) = E_{j+1,i}, \label{eq:increase}\\
T_{j} T_{j-1} \cdots T_{i+1} \big(x_{i}^+\big) = (-1)^{j-i}E_{i,j+1},\qquad T_{j} T_{j-1} \cdots T_{i+1} \big(x_{i}^-\big) = (-1)^{j-i}E_{j+1,i}. \label{eq:decrease}
\end{gather}
\end{lem}
\begin{proof}The assertion is easily proved by induction.
\end{proof}

\begin{lem}We have
\begin{alignat}{3}
&T_i(E_{i+1,j})= E_{i,j} \qquad && \text{if $j+1 \leq i \leq N-1$ \ or \ $i+2 \leq j$}, & \label{eq:TE-1}\\
&T_i(E_{j,i+1})= E_{j,i} \qquad &&\text{if $1 \leq i \leq j-2$ \ or \ $j \leq i-1$}, & \label{eq:TE-2}\\
&T_i(E_{i,1})= -E_{i+1,1} \qquad &&\text{if $2 \leq i \leq N-1$}, & \label{eq:TE-3}\\
&T_i(E_{2,i})= -E_{2,i+1} \qquad &&\text{if $3 \leq i \leq N-1$}. & \label{eq:TE-4}
\end{alignat}
\end{lem}
\begin{proof}The identity (\ref{eq:TE-2}) is deduced from (\ref{eq:TE-1}) by applying $\omega$. We use (\ref{eq:increase}) to show the other identities as
\begin{gather*}
T_i(E_{i+1,j}) = \begin{cases}
T_i T_j T_{j+1} {\cdots} T_{i-1}\big(x_i^-\big) = T_j T_{j+1} {\cdots} T_{i-2} \big(x_{i-1}^-\big) = E_{i,j} &\text{if $j+1 \leq i \leq N-1$}, \\
T_i T_{i+1} T_{i+2} \cdots T_{j-1}\big(x_j^+\big) = E_{i,j} & \text{if $i+2 \leq j$},
\end{cases}
\\
T_i(E_{i,1}) = T_i T_{1} T_{2} \cdots T_{i-2}\big(x_{i-1}^-\big) = T_{1} T_{2} \cdots T_{i-2} T_i \big(x_{i-1}^-\big)= T_{1} T_{2} \cdots T_{i-2} \big(\big[x_{i-1}^-,x_i^-\big]\big)\\
\hphantom{T_i(E_{i,1})}{} = \big[T_{1} T_{2} \cdots T_{i-2} \big(x_{i-1}^-\big),x_i^-\big] = [E_{i,1}, E_{i+1,i}] = -E_{i+1,1},\\
T_i(E_{2,i}) = T_i T_{2} T_{3} \cdots T_{i-2}\big(x_{i-1}^+\big) = T_{2} T_{3} \cdots T_{i-2} T_i \big(x_{i-1}^+\big) = T_{2} T_{3} \cdots T_{i-2} \big(\big[x_{i}^+,x_{i-1}^+\big]\big)\\
\hphantom{T_i(E_{2,i})}{} = \big[x_i^+, T_{2} T_{3} \cdots T_{i-2} \big(x_{i-1}^+\big)\big] = [E_{i,i+1}, E_{2,i}] = -E_{2,i+1}.\tag*{\qed}
\end{gather*}\renewcommand{\qed}{}
\end{proof}

\subsection{Evaluation map}

The evaluation map for the affine Yangian $\affY$ was introduced by Guay~\cite{MR2323534}. It is an affine analog of the well-known evaluation map from $Y(\fraksl_N)$ to $U(\gl_N)$. Let $U\big(\hat{\gl}_N\big)_{{\rm comp},-}$ be the completion of ${U}\big(\hat{\gl}_N\big)$ defined in \cite[Definition~2.5]{kodera_evaluation}. From now on, we regard the central element~$c$ of~$\hat{\gl}_N$ as a complex number.

The following result can be deduced from a formula for $\ev(H_{i,1})$ $(i \neq 0)$ where $H_{i,1}=h_{i,1} + (i/2)(\ve_1-\ve_2)h_i$, given in \cite[Section~6, pp.~462--463]{MR2323534}. See \cite{kodera_evaluation} for details.
\begin{thm}[\cite{MR2323534, kodera_evaluation}]Assume $\hbar c = N \ve_2$. Then there exists an algebra homomorphism $\ev \colon \affY \to U\big(\hat{\gl}_N\big)_{{\rm comp},-}$ uniquely determined by
\begin{gather*}
\ev(x_{i,0}^{+}) = x_{i}^{+}, \qquad \ev(x_{i,0}^{-}) = x_{i}^{-},\qquad \ev(h_{i,0}) = h_{i},\\
\ev(x_{i,1}^{+}) = \begin{cases}
(1+ N \varepsilon_2) x_{0}^{+} + \hbar \displaystyle\sum_{s \geq 0} \sum_{k=1}^N E_{k,1}(s+1) E_{N,k}(-s) \quad \text{if $i = 0$},\\
(1+ i \varepsilon_2) x_{i}^{+} + \hbar \displaystyle\sum_{s \geq 0} \bigg( \sum_{k=1}^i E_{k,i+1}(s) E_{i,k}(-s)\\
\qquad{} + \displaystyle \sum_{k=i+1}^N E_{k,i+1}(s+1) E_{i,k}(-s-1) \bigg) \quad \text{if $i \neq 0$},
\end{cases}\\
\ev(x_{i,1}^-) = \begin{cases}
(1+ N \varepsilon_2) x_{0}^{-} + \hbar \displaystyle\sum_{s \geq 0} \sum_{k=1}^N E_{k,N}(s) E_{1,k}(-s-1)\quad \text{if $i = 0$},\\
(1+ i \varepsilon_2) x_{i}^{-} + \hbar \displaystyle\sum_{s \geq 0} \bigg( \sum_{k=1}^i E_{k,i}(s) E_{i+1,k}(-s) \\
\qquad {} + \displaystyle\sum_{k=i+1}^N E_{k,i}(s+1) E_{i+1,k}(-s-1) \bigg) \quad \text{if $i \neq 0$},
\end{cases}\\
\ev(h_{i,1}) = \begin{cases}
(1+ N \varepsilon_2) h_{0} - \hbar E_{N,N} (E_{1,1}-c) \\
\qquad {}+ \hbar \displaystyle\sum_{s \geq 0} \sum_{k=1}^{N} ( E_{k,N}(s) E_{N,k}(-s) - E_{k,1}(s+1) E_{1,k}(-s-1) ) \quad \text{if $i = 0$},\\
(1+ i \varepsilon_2) h_{i} - \hbar E_{i,i}E_{i+1,i+1}+ \hbar \displaystyle\sum_{s \geq 0} \bigg( \sum_{k=1}^{i} E_{k,i}(s) E_{i,k}(-s) \\
\qquad{}+ \displaystyle\sum_{k=i+1}^{N} E_{k,i}(s+1) E_{i,k}(-s-1) - \displaystyle\sum_{k=1}^{i}E_{k,i+1}(s) E_{i+1,k}(-s) \\
\qquad {} - \displaystyle\sum_{k=i+1}^{N} E_{k,i+1}(s+1) E_{i+1,k}(-s-1) \bigg)\quad \text{if $i \neq 0$}.
\end{cases}
\end{gather*}
\end{thm}
We will use the following property.
\begin{prop}\label{prop:ev_and_auto}We have
\begin{gather*}
\mathrm{(i)}\ \omega \circ \ev = \ev \circ \,\omega, \qquad \mathrm{(ii)}\ \rho \circ \ev = \ev \circ\, \rho,\qquad \mathrm{(iii)}\ T_i \circ \ev = \ev \circ\, T_i.
\end{gather*}
\end{prop}
\begin{proof}The assertion (i) is immediate from the definition of $\ev$. The assertion (ii) is stated in \cite[Section~6, p.~463]{MR2323534} and a proof is given in \cite[Proposition~3.6]{kodera_evaluation}. Since~$\ev$ is the identity on the subalgebra $U\big(\affsl\big)$ and $T_i$ is defined via the generators of $U\big(\affsl\big)$, the assertion (iii) holds.
\end{proof}

\subsection{Construction of Heisenberg generators}

We construct elements $a_m$ ($m \in \bbZ$) of the affine Yangian $\affY$ such that $\ev(a_m) = \bfid(m)$ under the assumption $\ve_2 \neq 0$.

First consider the case $m=0$.
\begin{prop}\label{prop:zero-mode} We have
\begin{gather*} \ev \left( \sum_{i=0}^{N-1} \tilde{h}_{i,1} \right)= \ve_2 \bfid + c + \tfrac{\hbar}{2} c^2.
\end{gather*}
\end{prop}
\begin{proof}
Put
\begin{gather*}
A_0 = \sum_{s \geq 0} \sum_{k=1}^N E_{k,N}(s)E_{N,k}(-s),\qquad
B_0 = \sum_{s \geq 0} \sum_{k=1}^N E_{k,1}(s+1)E_{1,k}(-s-1),\\
A_i = \sum_{s \geq 0} \left( \sum_{k=1}^i E_{k,i}(s)E_{i,k}(-s) + \sum_{k=i+1}^N E_{k,i}(s+1)E_{i,k}(-s-1) \right), \qquad i \neq 0,\\
B_i = \sum_{s \geq 0} \left( \sum_{k=1}^i E_{k,i+1}(s)E_{i+1,k}(-s) + \sum_{k=i+1}^N E_{k,i+1}(s+1)E_{i+1,k}(-s-1) \right), \qquad i \neq 0,
\end{gather*}
so that
\begin{gather*}
\ev\big(\tilde{h}_{i,1}\big) = \begin{cases}
(1+ N \varepsilon_2) h_{0} - \tfrac{\hbar}{2} \big(E_{N,N}^2 + (E_{1,1}-c)^2\big) + \hbar(A_0-B_0)& \text{if $i=0$},\\
(1+ i \varepsilon_2) h_{i} - \tfrac{\hbar}{2} \big(E_{i,i}^2 + E_{i+1,i+1}^2\big) + \hbar(A_i-B_i)& \text{otherwise}.
\end{cases}
\end{gather*}
Here we use
\begin{gather*}
-\hbar E_{N,N}(E_{1,1}-c) - \tfrac{\hbar}{2} h_0^2 = - \tfrac{\hbar}{2} \big(E_{N,N}^2 + (E_{1,1}-c)^2\big),\\
-\hbar E_{i,i} E_{i+1,i+1} - \tfrac{\hbar}{2} h_i^2 = - \tfrac{\hbar}{2} \big(E_{i,i}^2 + E_{i+1,i+1}^2\big), \qquad i \neq 0.
\end{gather*}
Then we have $A_{i+1}-B_{i}=E_{i+1,i+1}^2$ for $i=0,\ldots,N-1$, where we regard $A_N=A_0$.
Therefore
\begin{gather*}
\ev \left( \sum_{i=0}^{N-1} \tilde{h}_{i,1} \right)= c + \ve_2 \left(N h_0 + \sum_{i=1}^{N-1} i h_i\right) -\tfrac{\hbar}{2}\left( 2\sum_{i=1}^N E_{i,i}^2 - 2cE_{1,1} + c^2\right) + \hbar \sum_{i=1}^N E_{i,i}^2 \\
\hphantom{\ev \left( \sum_{i=0}^{N-1} \tilde{h}_{i,1} \right)}{} = c + \ve_2 \sum_{i=1}^{N-1} i h_i + N \ve_2 h_0 + \hbar c E_{1,1} - \tfrac{\hbar}{2}c^2.
\end{gather*}
By the assumption $\hbar c = N\ve_2$, we have
\begin{gather*}
N \ve_2 h_0 + \hbar c E_{1,1} - \tfrac{\hbar}{2}c^2 = N \ve_2 (E_{N,N} + c) - \tfrac{\hbar}{2}c^2 = N \ve_2 E_{N,N} + \tfrac{\hbar}{2}c^2.
\end{gather*}
The assertion holds since we have $\bfid = \sum\limits_{i=1}^{N-1} i h_i + NE_{N,N}$.
\end{proof}

Assume $\hbar c = N \ve_2$ and $\ve_2 \neq 0$. Put
\begin{gather*}
a_0 = \dfrac{1}{\ve_2} \left( \sum_{i=0}^{N-1} \tilde{h}_{i,1} - c - \tfrac{\hbar}{2} c^2\right).
\end{gather*}
Then we have $\ev(a_0) = \bfid$ by Proposition~\ref{prop:zero-mode}.

Next consider the case $m \geq 1$. For each $i \in \bbZ/N\bbZ$ and a fixed $m$, define an element $w(i,m)$ of the affine Weyl group $\hat{W}$ by
\begin{gather*}
w(i,m) = t_{-\alpha_{i+1}}^{m-1} s_{i+1} s_{i+2} \cdots s_{i-3} s_{i-2}.
\end{gather*}
This element has the property $w(i,m)(\alpha_{i-1}) = -\alpha_{i} + m \delta$. Hence the elements
\begin{gather*}
\big[x_{i}^{+}, T_{w(i,m)}\big(x_{i-1}^{+}\big)\big],\qquad \big[x_{i}^{+}, T_{w(i,m)}\big(x_{i-1,1}^{+}\big)\big]
\end{gather*} have weight $m\delta$.
We will see
\begin{gather*}
\big[x_{i}^+, T_{w(i,m)}\big(x_{i-1}^+\big)\big] = (-1)^{m-1} \times \begin{cases}
h_{-\theta}(m) & \text{if $i=0$},\\
h_i(m) & \text{otherwise}
\end{cases}
\end{gather*}
in Lemma~\ref{lem:braket_pre}, and will compute the value of
\begin{gather*}
\ev \big( \big[x_{i}^{+}, T_{w(i,m)}\big(x_{i-1,1}^{+}\big)\big]\big) = \big[x_{i}^{+}, T_{w(i,m)}\ev\big(x_{i-1,1}^{+}\big)\big]
\end{gather*}
in Proposition~\ref{prop:braket}. Then we will take the summation over $i$ in Proposition~\ref{prop:Heisenberg}. The result will yield a construction of the elements $a_m$ in Theorem~\ref{thm:image}.

By Lemma~\ref{lem:rho_omega_and_T}(ii), $\rho \circ T_{w(i,m)} = T_{w(i-1,m)} \circ \rho$ holds. The case $i=1$ will be important. Note that $w(1,1)=s_2 s_3 \cdots s_{N-1}$ and $t_{-\alpha_2} = s_2 s_3 \cdots s_0 s_1 s_0 \cdots s_3$.
\begin{lem}\label{lem:diagonal2}We have
\begin{gather*}
T_{t_{-\alpha_2}}(E_{i,i}(s)) = \begin{cases}
E_{1,1}(s) & \text{if $i=1$}, \\
E_{2,2}(s) + \delta_{s,0}c & \text{if $i=2$}, \\
E_{3,3}(s) - \delta_{s,0}c & \text{if $i=3$}, \\
E_{i,i}(s) & \text{if $4 \leq i \leq N$}
\end{cases}
\end{gather*}
for any $s \in \bbZ$.
\end{lem}
\begin{proof}The assertion is easily proved by Lemma~\ref{lem:diagonal}.
\end{proof}
\begin{lem}\label{lem:i=1_partial}We have
\begin{gather*}
T_{w(1,m)}(E_{1,1}(s)) = E_{1,1}(s),\qquad T_{w(1,m)}(E_{N,N}(-s)) = E_{2,2}(-s) + \delta_{s,0}(m-1)c
\end{gather*}
for any $s \in \bbZ$.
\end{lem}
\begin{proof}The assertion follows from
\begin{gather*}
T_2 T_3 \cdots T_{N-1} (E_{1,1}(s)) = E_{1,1}(s), \qquad T_2 T_3 \cdots T_{N-1} (E_{N,N}(-s)) = E_{2,2}(-s)
\end{gather*}
and Lemma~\ref{lem:diagonal2}.
\end{proof}

\begin{lem}\label{lem:classical_part} We have
\begin{gather*} T_{w(i,m)}(x_{i-1}^+) = (-1)^{m-1} \times \begin{cases}
E_{1,N}(m-1) & \text{if $i=0$},\\
E_{i+1,i}(m) & \text{otherwise}.
\end{cases}
\end{gather*}
\end{lem}
\begin{proof}We prove the assertion for $i=1$. The other cases are deduced from this case by applying~$\rho$ and Lemma~\ref{lem:rho}. First assume $m=1$. Recall $T_{w(1,1)}=T_2 T_3 \cdots T_{N-1}$. We have
\begin{gather*}
T_2 T_3 \cdots T_{N-1}\big(x_0^+\big) = T_2 T_3 \cdots T_{N-2} \big(\big[x_{N-1}^+,x_0^+\big]\big) = \big[T_2 T_3 \cdots T_{N-2} \big(x_{N-1}^+\big),x_0^+\big] \\
\hphantom{T_2 T_3 \cdots T_{N-1}\big(x_0^+\big)}{} \overset{\text{by (\ref{eq:increase})}}{=} [E_{2,N}, E_{N,1}(1)] = E_{2,1}(1),
\end{gather*}
hence the case $m=1$ is proved. Next consider the case $m \geq 2$. Since the case $m=2$ yields the equality for general $m \geq 2$ inductively, it is enough to prove
\begin{gather*}
T_{t_{-\alpha_2}}(E_{2,1}(1))=-E_{2,1}(2).
\end{gather*}
Since $E_{2,1}(1)$ is invariant under $T_i$ for $3 \leq i \leq N-1$, we have
\begin{gather}
T_{t_{-\alpha_2}}(E_{2,1}(1)) = T_2 T_3 \cdots T_0 T_1 T_0 (E_{2,1}(1)).\label{eq:E21}
\end{gather}
We have
\begin{gather*}
T_0 T_1 T_0 (E_{2,1}(1)) = T_0 T_1 T_0 \big(T_2 T_3 \cdots T_{N-1}\big(x_0^+\big)\big) = T_0 T_1 T_2 T_3 \cdots T_{N-2}\big(x_{N-1}^+\big) = -E_{N,1}(2).
\end{gather*}
Here the second equality follows from the braid relations and Proposition~\ref{prop:iji_pre}, and the last from Lemma~\ref{lem:0_to_k}. Then the right-hand side of (\ref{eq:E21}) is
\begin{gather*}
T_2 T_3 \cdots T_{N-1} (-E_{N,1}(2)) \overset{\text{by Lemma~\ref{lem:higher}}}{=} - (T_2 T_3 \cdots T_{N-1} (E_{N,1}) )(2) \overset{\text{by (\ref{eq:TE-1})}}{=} -E_{2,1}(2),
\end{gather*}
hence the assertion is proved.
\end{proof}

\begin{lem}\label{lem:braket_pre}
We have
\begin{gather*} \big[x_{i}^+, T_{w(i,m)}\big(x_{i-1}^+\big)\big] = (-1)^{m-1} \times \begin{cases}
h_{-\theta}(m) & \text{if $i=0$},\\
h_i(m) & \text{otherwise}.
\end{cases}
\end{gather*}
\end{lem}

\begin{proof}
Immediate from Lemma~\ref{lem:classical_part}.
\end{proof}

\begin{lem}\label{lem:i=1}We have
\begin{gather*}
T_{w(1,m)}(E_{k,1}(s+1)E_{N,k}(-s)) \\
\qquad{} = (-1)^{m-1} \times\begin{cases}
E_{1,1}(s+1)E_{2,1}(-s+m-1) & \text{if $k=1$},\\
E_{3,1}(s-m+2)E_{2,3}(-s+2m-2) & \text{if $k=2$},\\
E_{k+1,1}(s+1)E_{2,k+1}(-s+m-1) & \text{if $3 \leq k \leq N-1$},\\
E_{2,1}(s+m)\Big( E_{2,2}(-s) + \delta_{s,0}(m-1)c \Big) & \text{if $k=N$}.
\end{cases}
\end{gather*}
\end{lem}

\begin{proof}We prove
\begin{gather}
T_{w(1,m)}(E_{k,1}(s+1)) = \begin{cases}
E_{1,1}(s+1) & \text{if $k=1$},\\
(-1)^{(m-1)(N+1)} (-E_{3,1}(s-m+2)) & \text{if $k=2$},\\
-E_{k+1,1}(s+1) & \text{if $3 \leq k \leq N-1$},\\
(-1)^{m-1} E_{2,1}(s+m) & \text{if $k=N$}
\end{cases}\label{eq:i=1former}
\end{gather}
and
\begin{gather}
T_{w(1,m)}(E_{N,k}(-s)) = \begin{cases}
(-1)^{m-1} E_{2,1}(-s+m-1) & \text{if $k=1$},\\
(-1)^{(m-1)N} (- E_{2,3}(-s+2m-2)) & \text{if $k=2$},\\
(-1)^{m-1}(- E_{2,k+1}(-s+m-1)) & \text{if $3 \leq k \leq N-1$},\\
E_{2,2}(-s) + \delta_{s,0}(m-1)c & \text{if $k=N$}.
\end{cases}\label{eq:i=1latter}
\end{gather}
The equalities (\ref{eq:i=1former}) for $k=1$ and (\ref{eq:i=1latter}) for $k=N$ follow from Lemma~\ref{lem:i=1_partial}.

Consider (\ref{eq:i=1former}) for $k=N$ and (\ref{eq:i=1latter}) for $k=1$. Note that Lemma~\ref{lem:classical_part} for $i=1$ is nothing but (\ref{eq:i=1former}) for $k=N$ and $s=0$. We can prove the other cases by applying $[-, E_{1,1}(\pm s)]$ to this case as
\begin{gather*}
\big[T_{w(1,m)}(E_{N,1}(1)), E_{1,1}(\pm s)\big] = T_{w(1,m)}\big(\big[E_{N,1}(1), E_{1,1}(\pm s)\big]\big) = T_{w(1,m)}\big(E_{N,1}(1 \pm s)\big).
\end{gather*}
Here we use the fact that $E_{1,1}(\pm s)$ is invariant under $T_{w(1,m)}$ proved in Lemma~\ref{lem:i=1_partial}. In the sequel, we prove (\ref{eq:i=1former}) and (\ref{eq:i=1latter}) for $2 \leq k \leq N-1$.

We prove (\ref{eq:i=1former}). First we consider the case $m=1$. We prove
\begin{gather*}
T_{w(1,1)} \big(E_{k,1}(s+1)\big)=-E_{k+1,1}(s+1)
\end{gather*}
for $2 \leq k \leq N-1$ by induction on $k$. Since $T_{w(1,1)}=T_2 T_3 \cdots T_{N-1}$ does not involve $T_0$, it is enough to prove $T_{w(1,1)}(E_{k,1})=-E_{k+1,1}$ by Lemma~\ref{lem:higher}. When $k=2$, we have
\begin{gather*}
T_2 T_3 \cdots T_{N-1}\big(x_1^-\big) = T_2\big(x_1^-\big) = -E_{3,1}.
\end{gather*}
Assume that the assertion holds for $k$. Since we have $E_{k+1,1}=-T_k(E_{k,1})$ by (\ref{eq:TE-3}),
\begin{gather*}
T_2 T_3 \cdots T_{N-1}(E_{k+1,1}) = -T_2 T_3 \cdots T_{N-1} T_k(E_{k,1}) = -T_{k+1} T_2 T_3 \cdots T_{N-1} (E_{k,1})\\
\hphantom{T_2 T_3 \cdots T_{N-1}(E_{k+1,1})}{} = T_{k+1}(E_{k+1,1}) \overset{\text{by (\ref{eq:TE-3})}}{=} -E_{k+2,1}.
\end{gather*}
Thus we have proved the case $m=1$. Next we consider the case $m \geq 2$. Since the case $m=2$ yields the equality for general $m \geq 2$ inductively, it is enough to prove:
\begin{gather}
T_{t_{-\alpha_2}}\big(E_{3,1}(s+1)\big) = (-1)^{N+1}E_{3,1}(s), \label{eq:m=2_k=2}\\
T_{t_{-\alpha_2}}\big(E_{k+1,1}(s+1)\big) = E_{k+1,1}(s+1) \qquad \text{if $3 \leq k \leq N-1$}. \label{eq:m=2_k+1}
\end{gather}
If we prove the assertion for $s=0$, we can prove the other cases by applying $[-,E_{1,1}(s)]$ to this case by using the fact that $E_{1,1}(s)$ is invariant under $T_{t_{-\alpha_2}}$ proved in Lemma~\ref{lem:diagonal2}.
We prove~(\ref{eq:m=2_k=2}) for $s=0$. Since we have $E_{3,1}=-T_2\big(x_1^-\big)$, the left-hand side of~(\ref{eq:m=2_k=2}) for $s=0$ is
\begin{gather*}
-T_2 T_3 \cdots T_0 T_1 T_0 \cdots T_3 \big(T_2 \big(x_1^-\big)(1)\big) \overset{\text{by (\ref{eq:decrease})}}{=} -T_2 T_3 \cdots T_0 T_1 T_0\big( (-1)^{N-2} E_{N,1}(1) \big) \\
\qquad {} = (-1)^{N-1} T_2 T_3 \cdots T_0 T_1 T_0 \big(x_0^+\big) = (-1)^N T_2 T_3 \cdots T_{N-1} \big(x_1^-\big) \\
\qquad{} = (-1)^N T_2\big(x_1^-\big) = (-1)^{N+1} E_{3,1}.
\end{gather*}
We prove (\ref{eq:m=2_k+1}) for $s=0$ by backward induction on $k$. The case $k=N-1$ is proved as
\begin{gather*}
T_2 T_3 \cdots T_0 T_1 T_0 \cdots T_3\big(x_0^+\big) = T_2 T_3 \cdots T_0 T_1 \big(x_{N-1}^+\big) = T_2 T_3 \cdots T_{N-2} \big(x_0^+\big) = x_0^+.
\end{gather*}
Assume that the assertion holds for $k$. Since we have $E_{k,1} = T_k(E_{k+1,1})$ by~(\ref{eq:TE-1}),
\begin{gather*}
T_{t_{-\alpha_2}}\big(E_{k,1}(1)\big) = T_{t_{-\alpha_2}} T_{k}\big(E_{k+1,1}(1)\big) = T_{k} T_{t_{-\alpha_2}}\big(E_{k+1,1}(1)\big) = T_k \big(E_{k+1,1}(1)\big) = E_{k,1}(1).
\end{gather*}
Here the second equality holds by $k \geq 3$. Thus we have proved (\ref{eq:i=1former}).

We prove (\ref{eq:i=1latter}). First we consider the case $m=1$. We prove
\begin{gather*}
T_{w(1,1)} \big(E_{N,k}(-s)\big)=-E_{N,k+1}(-s)
\end{gather*}
for $2 \leq k \leq N-1$ by backward induction on $k$. By Lemma~\ref{lem:higher}, it is enough to prove $T_{w(1,1)}(E_{N,k})=-E_{2,k+1}$. When $k=N-1$, we have
\begin{gather*}
T_2 T_3 \cdots T_{N-1}\big(x_{N-1}^-\big) = -T_2 T_3 \cdots T_{N-2}\big(x_{N-1}^+\big) = -E_{2,N}
\end{gather*}
by (\ref{eq:increase}). Assume that the assertion holds for $k$. Since we have $E_{N,k-1}=T_{k-1}(E_{N,k})$ by (\ref{eq:TE-2}),
\begin{gather*}
T_2 T_3 \cdots T_{N-1}\big(E_{N,k-1}\big) = T_2 T_3 \cdots T_{N-1} T_{k-1}\big(E_{N,k}\big) = T_k T_2 T_3 \cdots T_{N-1} \big(E_{N,k}\big)\\
\hphantom{T_2 T_3 \cdots T_{N-1}\big(E_{N,k-1}\big)}{} = -T_k(E_{2,k+1}) \overset{\text{by (\ref{eq:TE-2})}}{=} -E_{2,k}.
\end{gather*}
Thus we have proved the case $m=1$. Next we consider the case $m \geq 2$. Since the case $m=2$ yields the equality for general $m \geq 2$ inductively, it is enough to prove:
\begin{gather}
T_{t_{-\alpha_2}}\big(E_{2,3}(-s)\big) = (-1)^N E_{2,3}(-s+2), \label{eq:m=2_k=2_second}\\
T_{t_{-\alpha_2}}\big(E_{2,k+1}(-s)\big) = -E_{2,k+1}(-s+1) \qquad \text{if $3 \leq k \leq N-1$}. \label{eq:m=2_k+1_second}
\end{gather}
If we prove the assertion for $s=0$, we can prove the other cases by applying $[E_{2,2}(-s),-]$ to this case by using the fact that $E_{2,2}(-s)$ is invariant under $T_{t_{-\alpha_2}}$ for $s \geq 1$ proved in Lemma~\ref{lem:diagonal2}. We prove (\ref{eq:m=2_k=2_second}) for $s=0$. Since we have
\begin{gather*}
T_{N-1} T_{N-2} \cdots T_3 \big(x_2^+\big) = (-1)^{N-3} E_{2,N}
\end{gather*}
by (\ref{eq:decrease}), and
\begin{gather*}
T_0 T_1 T_0 \big(E_{2,N}\big) = T_1 T_0 T_1 \big(E_{2,N}\big) \overset{\text{by (\ref{eq:TE-1})}}{=} T_1 T_0 \big(E_{1,N}\big)
\overset{\text{by Lemma~\ref{lem:0_to_k}}}{=} T_1\big({-}E_{N,1}(2)\big),
\end{gather*}
the left-hand side of (\ref{eq:m=2_k=2_second}) for $s=0$ is
\begin{gather*}
(-1)^{N-2} T_2 T_3 \cdots T_{N-1} T_1 \big(E_{N,1}(2)\big) = (-1)^{N-2} T_2 T_1 T_3 \cdots T_{N-1} \big(E_{N,1}(2)\big) \\
\qquad{} \overset{\text{by (\ref{eq:TE-1})}}{=} (-1)^{N-2} T_2 T_1 \big(E_{3,1}\big) (2)= (-1)^{N-2} T_2 T_1 \big({-}T_2\big(x_1^-\big)\big)(2) \\
\qquad {} = (-1)^{N-1} T_2 \big(x_2^-\big)(2) = (-1)^N E_{2,3}(2).
\end{gather*}
We prove (\ref{eq:m=2_k+1_second}) for $s=0$ by induction on $k$. Assume $k=3$. Since we have $E_{2,4}=T_2\big(x_3^+\big)$, the left-hand side of (\ref{eq:m=2_k+1_second}) for $s=0$ is
\begin{gather*}
T_2 T_3 \cdots T_0 T_1 T_0 \cdots T_3\big(T_2\big(x_3^+\big)\big) = T_2 T_3 \cdots T_0 T_1 T_0 \cdots T_{4} \big(x_{2}^+\big) = T_2 T_3 \cdots T_0 T_1 \big(x_2^+\big) \\
\qquad \overset{\text{by Lemma~\ref{lem:0_to_k}}}{=} T_2 T_3 \cdots T_{N-1} \big(E_{N,3}(1)\big)\overset{\text{by (\ref{eq:TE-1})} }{=} T_2 T_3 (E_{4,3}) (1) = - T_2\big(x_3^+\big)(1) = -E_{2,4}(1).
\end{gather*}
Assume that the assertion holds for $k$.
Since we have $E_{2,k+2} = -T_{k+1}(E_{2,k+1})$ by (\ref{eq:TE-4}), the left-hand side is
\begin{gather*}
T_{t_{-\alpha_2}}\big(E_{2,k+2}\big) = -T_{t_{-\alpha_2}} T_{k+1}\big(E_{2,k+1}\big) = -T_{k+1} T_{t_{-\alpha_2}}\big(E_{2,k+1}\big)\\
\hphantom{T_{t_{-\alpha_2}}\big(E_{2,k+2}\big)}{} = T_{k+1} \big(E_{2,k+1}\big)(1) = -E_{2,k+2}(1).
\end{gather*}
Here the second equality holds by $k \geq 3$.
\end{proof}

For a fixed $m \geq 1$, put
\begin{gather*}
S_{i,j}(p) = E_{i,j}(p)E_{j,i}(m-p).
\end{gather*}

\begin{prop}\label{prop:braket}
We have
\begin{gather}
\big[x_i^+, T_{w(i,m)} \ev\big(x_{i-1,1}^+\big)\big] \nonumber\\
= (-1)^{m-1} \left(A_i + \hbar \sum_{k=1}^N (P_{i,k}-Q_{i,k}) - \hbar \sum_{s=0}^{m-2} E_{i,i}(s+1)E_{i+1,i+1}(-s+m-1) \right),\label{eq:braket}
\end{gather}
where
\begin{gather*}
A_i = \begin{cases}
(1+(N-1)\ve_2)h_{-\theta}(m) + (m-1)\hbar ch_{-\theta}(m)+ \hbar c E_{N,N}(m) & \text{if $i=0$}, \\
(1+N\ve_2)h_1(m) + (m-1)\hbar ch_{1}(m)& \text{if $i=1$}, \\
(1+(i-1)\ve_2)h_{i}(m) + (m-1)\hbar ch_{i}(m) + \hbar c h_{i}(m) & \text{if $2 \leq i \leq N-1$},
\end{cases}
\\
P_{i,k} = \sum_{s \geq 0} S_{k,i}(p(i,k)), \qquad Q_{i,k} = \sum_{s \geq 0} S_{k,i+1}(q(i,k)).
\end{gather*}
The integers $p(i,k)$ and $q(i,k)$ are given by
\begin{gather*}
p(0,k)= \begin{cases}
s+m-1 & \text{if $k=1$},\\
s-m+1 & \text{if $k=2$},\\
s & \text{if $3 \leq k \leq N-1$},\\
s+1 & \text{if $k=N$},
\end{cases}\\
p(i,k)= \begin{cases}
s & \text{if $1 \leq k \leq i-1$},\\
s+1 & \text{if $k=i$},\\
s+m & \text{if $k=i+1$},\\
s-m+2 & \text{if $k=i+2$},\\
s+1 & \text{if $i+3 \leq k \leq N$},
\end{cases}\qquad 1 \leq i \leq N-2,\\
p(N-1,k)= \begin{cases}
s-m+1 & \text{if $k=1$},\\
s & \text{if $2 \leq k \leq N-2$},\\
s+1 & \text{if $k=N-1$},\\
s+m & \text{if $k=N$},
\end{cases}\qquad
q(i,k)=\begin{cases}
p(0,k)+1 & \text{if $i=0$},\\
p(i,k) &\text{otherwise}.
\end{cases}
\end{gather*}
\end{prop}

\begin{proof}We prove the assertion for $i=1$. By Lemma~\ref{lem:i=1}, we have
\begin{gather*}
(-1)^{m-1}T_{w(1,m)}\ev\big(x_{0,1}^+\big) \\
 \qquad{} = (-1)^{m-1}(1+N\ve_2)T_{w(1,m)}\big(x_0^+\big) + (-1)^{m-1}\hbar \sum_{s \geq 0} \sum_{k=1}^N T_{w(1,m)}\big(E_{k,1}(s+1)E_{N,k}(-s)\big) \\
\qquad{} =(-1)^{m-1}(1+N\ve_2)T_{w(1,m)}\big(x_0^+\big) + (m-1)\hbar c E_{2,1}(m)\\
\qquad \quad{} + \hbar \sum_{s \geq 0} \bigg( E_{1,1}(s+1)E_{2,1}(-s+m-1) + E_{2,1}(s+m)E_{2,2}(-s) \\
 \qquad \quad{}+ E_{3,1}(s-m+2)E_{2,3}(-s+2m-2) + \sum_{k=4}^N E_{k,1}(s+1)E_{2,k}(-s+m-1) \bigg).
\end{gather*}
Therefore
\begin{gather}
(-1)^{m-1} [x_1^+, T_{w(1,m)} \ev(x_{0,1}^+)] = (1+N\ve_2)h_1(m) + (m-1)\hbar c h_1(m) \nonumber\\
{}+ \hbar \sum_{s \geq 0} \bigg( {-}S_{1,2}(s+1) + E_{1,1}(s+1)h_1(-s+m-1) + h_1(s+m)E_{2,2}(-s) + S_{2,1}(s+m)\nonumber\\
{}-S_{3,2}(s-m+2) + S_{3,1}(s-m+2) + \sum_{k=4}^N \big( {-}S_{k,2}(s+1) + S_{k,1}(s+1) \big) \bigg). \label{eq:braket_i=1}
\end{gather}
Here we use Lemma~\ref{lem:braket_pre}. Since we have
\begin{gather*}
\sum_{s \geq 0} \big( E_{1,1}(s+1)h_1(-s+m-1) + h_1(s+m)E_{2,2}(-s) \big)\\
\qquad {}=\sum_{s \geq 0} \big( S_{1,1}(s+1)-S_{2,2}(s+m) \big) - \sum_{s=0}^{m-2} E_{1,1}(s+1)E_{2,2}(-s+m-1),
\end{gather*}
the right-hand side of (\ref{eq:braket_i=1}) is equal to
\begin{gather*}
(1+N\ve_2)h_1(m) + (m-1)\hbar c h_1(m) \\
\qquad{} + \hbar \sum_{s \geq 0} \bigg( \bigg( S_{1,1}(s+1) + S_{2,1}(s+m) + S_{3,1}(s-m+2) + \sum_{k=4}^N S_{k,1}(s+1) \bigg) \\
\qquad{} - \bigg( S_{1,2}(s+1) + S_{2,2}(s+m) + S_{3,2}(s-m+2) + \sum_{k=4}^N S_{k,2}(s+1) \bigg) \bigg)\\
\qquad{} - \hbar \sum_{s=0}^{m-2} E_{1,1}(s+1)E_{2,2}(-s+m-1).
\end{gather*}
Hence the assertion holds for $i=1$. Then apply $\rho$ to (\ref{eq:braket}) for $i=1$. The left-hand side is
\begin{gather*}
\begin{split}&
 \rho\big(\big[x_1^+, T_{w(1,m)} \ev\big(x_{0,1}^+\big)\big]\big)=\big[x_0^+, T_{w(0,m)} \ev\big(x_{0,1}^++\ve_2 x_0^+\big)\big]\\
& \hphantom{\rho\big(\big[x_1^+, T_{w(1,m)} \ev\big(x_{0,1}^+\big)\big]\big)}{} = \big[x_0^+, T_{w(0,m)} \ev\big(x_{0,1}^+\big)\big] + \ve_2 h_{-\theta}(m).
\end{split}
\end{gather*}
Here we use Proposition~\ref{prop:ev_and_auto}(ii) and Lemma~\ref{lem:braket_pre}. Consider the right-hand side. We can see
\begin{gather*}
\rho(A_1) = A_0 + \ve_2 h_{-\theta}(m) - \hbar c E_{N,N}(m),\\
\rho(P_{1,k}) = \begin{cases}
P_{0,N} + E_{N,N}(m)c & \text{if $k=1$},\\
P_{0,k-1} & \text{if $2 \leq k \leq N$},
\end{cases}\qquad
\rho(Q_{1,k}) = Q_{0,k-1},\\
\rho\big(E_{1,1}(s+1)E_{2,2}(-s+m-1)\big) = E_{N,N}(s+1)E_{1,1}(-s+m-1).
\end{gather*}
Hence the assertion holds for $i=0$. Then apply $\rho$ to (\ref{eq:braket}) for $i=0$. Similarly the left-hand side is
\begin{gather*}
 \rho\big(\big[x_0^+, T_{w(0,m)} \ev\big(x_{N-1,1}^+\big)\big]\big)= \big[x_{N-1}^+, T_{w(N-1,m)} \ev\big(x_{N-2,1}^+\big)\big] + \ve_2 h_{N-1}(m),
\end{gather*}
and for the right-hand side, we can see
\begin{gather*}
\rho(A_0) = A_{N-1} + \ve_2 h_{N-1}(m) + \hbar c E_{N,N}(m),\\
\rho(P_{0,k}) = P_{N-1,k-1}, \qquad \rho(Q_{0,k}) = \begin{cases}
Q_{N-1,N} - E_{N,N}(m)c & \text{if $k=1$},\\
Q_{N-1,k-1} & \text{if $2 \leq k \leq N$},
\end{cases}\\
\rho\big(E_{N,N}(s+1)E_{1,1}(-s+m-1)\big) = E_{N-1,N-1}(s+1)E_{N,N}(-s+m-1).
\end{gather*}
Note that $c$ never appears in the last equality as $-s+m-1$ cannot be $0$ for $0 \leq s \leq m-2$. Hence the assertion holds for $i=N-1$. Continuing this process we prove the assertions for $i=N-2, N-3,\ldots,2$ since we have
\begin{gather*}
\rho(A_{i+1}) = A_{i} + \ve_2 h_{i}(m),\qquad \rho(P_{i+1,k}) = P_{i,k-1}, \qquad \rho(Q_{i+1,k}) =Q_{i,k-1},\\
\rho\big(E_{i+1,i+1}(s+1)E_{i+2,i+2}(-s+m-1)\big) = E_{i,i}(s+1)E_{i+1,i+1}(-s+m-1).\tag*{\qed}
\end{gather*}\renewcommand{\qed}{}
\end{proof}
\begin{prop}\label{prop:Heisenberg}
We have
\begin{gather}
\ev\left( (-1)^{m-1} \sum_{i=0}^{N-1} \big[x_{i}^+, T_{w(i,m)}\big(x_{i-1,1}^+\big)\big] \right) = \ve_2 \bfid(m) + \hbar R_m, \label{eq:Heisenberg}
\end{gather}
where
\begin{gather*}
R_m= (-1)^{m-1}\sum_{i=0}^{N-1} x_{i}^+ T_{w(i,m)}\big(x_{i-1}^+\big) + \sum_{p=1}^{m-1} \sum_{1 \leq i \leq j \leq N-1} h_i(p)h_j(m-p) \\
\hphantom{R_m=}{} - \sum_{p=-m+2}^{m-1} \bigg( E_{1,N}(p-1)E_{N,1}(m-p+1) + \sum_{i=0}^{N-2} E_{i+2,i+1}(p)E_{i+1,i+2}(m-p) \bigg) \\
\hphantom{R_m=}{} + \sum_{p=-m+2}^{0} \bigg( E_{1,N-1}(p-1)E_{N-1,1}(m-p+1) +E_{2,N}(p-1)E_{N,2}(m-p+1) \\
\hphantom{R_m=}{} + \sum_{i=0}^{N-3} E_{i+3,i+1}(p)E_{i+1,i+3}(m-p) \bigg).
\end{gather*}
\end{prop}

\begin{rem}The point of the statement of Proposition~\ref{prop:Heisenberg} is as follows: although each term $\ev\big(\big[x_{i}^+, T_{w(i,m)}\big(x_{i-1,1}^+\big)\big]\big)$ lies in the completion of $U\big(\hat{\gl}_N\big)$, we obtain $R_m$ an element of $U\big(\affsl\big)$ as a remainder term after cancellation.
\end{rem}

\begin{proof} We use the notation in Proposition~\ref{prop:braket}. We have
\begin{gather*}
\sum_{i=0}^{N-1} A_i = \ve_2 \left((N-1)h_{-\theta} + N h_1 + \sum_{i=2}^{N-1}(i-1)h_i\right)(m) + \hbar c \left(E_{N,N}+\sum_{i=2}^{N-1} h_i\right)(m) \\
\hphantom{\sum_{i=0}^{N-1} A_i}{} = \ve_2 \left( (N-1)h_{-\theta} + Nh_1 + \sum_{i=2}^{N-1} (N+i-1) h_i + NE_{N,N} \right)(m) \\
\hphantom{\sum_{i=0}^{N-1} A_i}{} = \ve_2 \left( \sum_{i=1}^{N-1} i h_i + NE_{N,N}\right)(m)=\ve_2 \bfid(m).
\end{gather*}
In the second equality we use the condition $\hbar c = N\ve_2$. We compute $\sum\limits_{k=1}^{N} (P_{i+1,k}-Q_{i,k})$ as follows:
\begin{gather*}
\sum_{k=1}^{N} (P_{0,k}-Q_{N-1,k}) = - \sum_{s=0}^{2m-3}S_{1,N}(s-m+1) + \sum_{s=0}^{m-2}S_{2,N}(s-m+1)\\
\hphantom{\sum_{k=1}^{N} (P_{0,k}-Q_{N-1,k})=}{} +S_{N-1,N}(0)+\sum_{s=0}^{m-2} S_{N,N}(s+1),
\\
\sum_{k=1}^{N} (P_{1,k}-Q_{0,k}) = \sum_{s=0}^{m-2} S_{1,1}(s+1) - \sum_{s=0}^{2m-3}S_{2,1}(s-m+2) \\
\hphantom{\sum_{k=1}^{N} (P_{1,k}-Q_{0,k}) =}{} + \sum_{s=0}^{m-2}S_{3,1}(s-m+2)+S_{N,1}(1),
\\
\sum_{k=1}^{N} (P_{i+1,k}-Q_{i,k}) = S_{i,i+1}(0) + \sum_{s=0}^{m-2} S_{i+1,i+1}(s+1)- \sum_{s=0}^{2m-3}S_{i+2,i+1}(s-m+2) \\
\hphantom{\sum_{k=1}^{N} (P_{i+1,k}-Q_{i,k}) =}{} + \sum_{s=0}^{m-2}S_{i+3,i+1}(s-m+2) \qquad \text{for}\ 1 \leq i \leq N-3,
\\
\sum_{k=1}^{N} (P_{N-1,k}-Q_{N-2,k})= \sum_{s=0}^{m-2}S_{1,N-1}(s-m+1)+S_{N-2,N-1}(0)\\
\hphantom{\sum_{k=1}^{N} (P_{N-1,k}-Q_{N-2,k})=}{} +\sum_{s=0}^{m-2} S_{N-1,N-1}(s+1) - \sum_{s=0}^{2m-3}S_{N,N-1}(s-m+2).
\end{gather*}
Hence the assertion holds by
\begin{gather*}
(-1)^{m-1} x_i^+ T_{w(i,m)}\big(x_{i-1}^+\big) = \begin{cases}
S_{N,1}(1) & \text{if $i=0$,}\\
S_{i,i+1}(0) & \text{otherwise}
\end{cases}
\end{gather*}
and
\begin{gather*}
\sum_{i=0}^{N-1} \big( S_{i,i}(s+1) - E_{i,i}(s+1) E_{i+1,i+1}(-s+m-1) \big) \\
\qquad{}= \left( \sum_{i=1}^{N-1} E_{i,i}(s+1) h_{i}(-s+m-1) \right)+ E_{N,N}(s+1)h_{-\theta}(-s+m-1)\\
\qquad{} = \sum_{i=1}^{N-1} (E_{i,i}-E_{N,N})(s+1) h_{i}(-s+m-1) = \sum_{i=1}^{N-1} \sum_{j \geq i} h_j(s+1) h_i(-s+m-1).\tag*{\qed}
\end{gather*}\renewcommand{\qed}{}
\end{proof}

Applying $\omega$ to (\ref{eq:Heisenberg}), we obtain
\begin{gather*}
\ev\left( (-1)^{m-1} \sum_{i=0}^{N-1} \big[T_{w(i,m)}\big(x_{i-1,1}^-\big), x_{i}^-\big] \right) = \ve_2 \bfid(-m) + \hbar \omega(R_m).
\end{gather*}
Now the following theorem has been proved.
\begin{thm}\label{thm:image}Assume $\hbar c =N\ve_2$ and $\ve_2 \neq 0$. Let $R_m$ $(m \geq 1)$ be the element of $U(\hat{\mathfrak{sl}}_N) \subset \affY$ as in Proposition~{\rm \ref{prop:Heisenberg}}, and define for $m \in \bbZ$,
\begin{gather*}
a_{m} = \dfrac{1}{\ve_2} \times \begin{cases}
\displaystyle \left( (-1)^{m-1} \sum_{i=0}^{N-1} \big[x_{i}^+, T_{w(i,m)}\big(x_{i-1,1}^+\big)\big] - \hbar R_m\right) & \text{if $m>0$},\\
\displaystyle \left( \sum_{i=0}^{N-1} \tilde{h}_{i,1} - c - \tfrac{\hbar}{2} c^2\right) & \text{if $m=0$},\\
\displaystyle \left((-1)^{m-1} \sum_{i=0}^{N-1} \big[T_{w(i,-m)}\big(x_{i-1,1}^-\big),x_{i}^-\big] - \hbar \omega(R_{-m})\right) & \text{if $m<0$}.
\end{cases}
\end{gather*}
Then we have $\ev(a_m) = \bfid(m)$. In particular, the image of the evaluation map $\ev$ contains $U\big(\hat{\gl}_N\big)$.
\end{thm}
